% ================================================================
%  CHAPTER 3 -- AMORTISATION & NEURAL BAYES ESTIMATORS
%  Status: DRAFT v1 (2026-05-18). Adapted from essay Sec 3-5.
% ================================================================
\chapter{Amortisation and Neural Bayes Estimators}
\label{ch:amortization}

MCMC solves each inference problem in isolation:
for every new dataset the chain must be run afresh, at a cost that is
incurred each time. This chapter
develops an alternative. Rather than approximating the full posterior
for one dataset at a time, we target specific \emph{posterior
functionals} (summaries of the form $\E[g(\bm{\theta}) \given \mathbf{y}]$) and we learn them with a single
estimator map $\hat{\bm{\theta}}$, trained once, that returns the desired summary for any
dataset drawn from the model. The idea is to exploit the \emph{universal
function-approximation} capability of neural networks to approximate
the optimal mapping from data to these summaries.

The functionals we target are not arbitrary. Each is the solution of a
decision problem: for a suitable loss function, the posterior summary
$\E[g(\bm{\theta}) \given \mathbf{y}]$ is the \emph{Bayes estimator},
the estimator that is optimal with respect to the posterior
distribution. Section~\ref{sec:bayes-estimator} develops this decision
theory and shows how the choice of loss selects which posterior
summary the Bayes estimator returns: the posterior mean under
squared-error loss, posterior quantiles under quantile loss.

The Bayes estimator is a fixed function of the data, but evaluating it
directly is generally as difficult as the original inference problem.
Section~\ref{sec:nbe} introduces \emph{amortised inference} (which
approximates this function once, up front, and reuses the
approximation across all subsequent datasets, so that the fitting cost
is amortised over every later inference) and makes it concrete: the
approximating function is a neural network which we call a \emph{neural
Bayes estimator}
%
\begin{equation}
  \hat{\bm{\theta}}_{\bm{\gamma}}(\mathbf{y})
\approx \E[g(\bm{\theta}) \given \mathbf{y}],
\end{equation}
%
trained on simulated data to minimise the empirical Bayes risk. We discuss how the network architecture should match
the structure of the data (DeepSets for exchangeable observations,
recurrent networks for sequential data, feed-forward networks
otherwise) and close Section~\ref{sec:normal-example} with a
worked example estimating the parameters of a normal model, where the
amortised estimator is checked against an MCMC reference.


% ----------------------------------------------------------------
\section{The Bayes estimator framework}
\label{sec:bayes-estimator}

Point estimation can be formalised through decision
theory~\citep{berger1985decision}. A point
estimator is a function $\hat{\bm{\theta}} : \mathcal{Y} \to \Theta$
mapping observed data $\mathbf{y}$ to an estimate
$\hat{\bm{\theta}}(\mathbf{y})$. A \emph{loss function}
$L(\bm{\theta}, \hat{\bm{\theta}})$ is a function that quantifies
the cost of reporting
$\hat{\bm{\theta}}$ when the true parameter is $\bm{\theta}$.

The quality of an estimator at a fixed parameter value is its expected
loss over data realisations, the \emph{frequentist risk}
%
\begin{equation}
  R(\bm{\theta}, \hat{\bm{\theta}})
  = \E_{\mathbf{y} \given \bm{\theta}}
    \big[L(\bm{\theta}, \hat{\bm{\theta}}(\mathbf{y}))\big].
  \label{eq:frequentist-risk}
\end{equation}
%
Averaging the frequentist risk over the prior gives the \emph{Bayes
risk},
%
\begin{equation}
  r(\hat{\bm{\theta}})
  = \E_{\bm{\theta} \sim p(\bm{\theta})}
    \big[R(\bm{\theta}, \hat{\bm{\theta}})\big]
  = \E_{\bm{\theta}, \mathbf{y}}
    \big[L(\bm{\theta}, \hat{\bm{\theta}}(\mathbf{y}))\big],
  \label{eq:bayes-risk}
\end{equation}
%
the second equality following because $p(\bm{\theta})\,
p(\mathbf{y} \given \bm{\theta})$ is the joint density of
$(\bm{\theta}, \mathbf{y})$.

\begin{definition}[Bayes estimator]
\label{def:bayes-estimator}
The \emph{Bayes estimator} is the estimator minimising the Bayes risk,
%
\begin{equation}
  \hat{\bm{\theta}}^{*} (\cdot)
  = \arg\min_{\hat{\bm{\theta}}}\;
    r(\hat{\bm{\theta}} (\cdot)).
  \label{eq:bayes-estimator}
\end{equation}
\end{definition}

By the law of iterated expectation the Bayes
risk~\eqref{eq:bayes-risk} can be rewritten as
$r(\hat{\bm{\theta}}) = \E_{\mathbf{y}}\big[\E_{\bm{\theta} \given
\mathbf{y}}[L(\bm{\theta}, \hat{\bm{\theta}}(\mathbf{y}))]\big]$.
Because the estimator is free to choose $\hat{\bm{\theta}}(\mathbf{y})$
separately for each $\mathbf{y}$, the Bayes risk is minimised by
minimising the inner \emph{posterior expected loss}
$\E_{\bm{\theta} \given \mathbf{y}}[L(\bm{\theta},
\hat{\bm{\theta}})]$ pointwise. The Bayes estimator~\eqref{eq:bayes-estimator} is therefore
%
\begin{equation}
  \hat{\bm{\theta}}^{*}(\mathbf{y})
  = \arg\min_{\hat{\bm{\theta}}}\;
    \E_{\bm{\theta} \given \mathbf{y}}
    \big[L(\bm{\theta}, \hat{\bm{\theta}})\big].
  \label{eq:pointwise}
\end{equation}

\subsection{Loss functions and posterior summaries}
\label{sec:loss-summaries}

A crucial feature of the Bayes estimator is that by choosing different loss functions, 
it can be made to target different summaries of the posterior.
The two results below show the loss functions we need to target the posterior mean and posterior quantiles.  

\begin{proposition}[Squared-error loss]
\label{prop:squared-error}
Under the squared-error loss
$L(\bm{\theta}, \hat{\bm{\theta}}) =
\|\bm{\theta} - \hat{\bm{\theta}}\|^2$, the Bayes estimator is the
posterior mean $\hat{\bm{\theta}}^{*}(\mathbf{y}) =
\E[\bm{\theta} \given \mathbf{y}]$.
\end{proposition}

\begin{proof}
Fix the data $\mathbf{y}$, so that $\hat{\bm{\theta}}$ is the constant
vector to be chosen, and expand the squared norm inside the posterior
expectation:
%
\begin{align}
  \E_{\bm{\theta} \given \mathbf{y}}
  \big[\|\bm{\theta} - \hat{\bm{\theta}}\|^2\big]
  &= \E_{\bm{\theta} \given \mathbf{y}}
     \big[(\bm{\theta} - \hat{\bm{\theta}})^{\!\top}
     (\bm{\theta} - \hat{\bm{\theta}})\big]
   \\[2pt]
  &= \E_{\bm{\theta} \given \mathbf{y}}
     \big[\bm{\theta}^{\!\top}\bm{\theta}
     - 2\,\hat{\bm{\theta}}^{\!\top}\bm{\theta}
     + \hat{\bm{\theta}}^{\!\top}\hat{\bm{\theta}}\big]
   \\[2pt]
  &= \E[\bm{\theta}^{\!\top}\bm{\theta} \given \mathbf{y}]
     - 2\,\hat{\bm{\theta}}^{\!\top}\,
       \E[\bm{\theta} \given \mathbf{y}]
     + \hat{\bm{\theta}}^{\!\top}\hat{\bm{\theta}},
\end{align}
%
where the last line uses linearity of expectation. Differentiating with respect to $\hat{\bm{\theta}}$ using the gradient identities
$\nabla_{\hat{\bm{\theta}}}(\hat{\bm{\theta}}^{\!\top}\mathbf{a})
= \mathbf{a}$ and
$\nabla_{\hat{\bm{\theta}}}
(\hat{\bm{\theta}}^{\!\top}\hat{\bm{\theta}}) = 2\,\hat{\bm{\theta}}$ gives
%
\begin{equation}
  \nabla_{\hat{\bm{\theta}}}\;
  \E_{\bm{\theta} \given \mathbf{y}}
  \big[\|\bm{\theta} - \hat{\bm{\theta}}\|^2\big]
  = -2\,\E[\bm{\theta} \given \mathbf{y}]
  + 2\,\hat{\bm{\theta}}.
\end{equation}
%
The objective is the sum of a constant, a term linear in
$\hat{\bm{\theta}}$, and the convex quadratic
$\hat{\bm{\theta}}^{\!\top}\hat{\bm{\theta}}$; it is therefore convex
in $\hat{\bm{\theta}}$, and its unique minimiser is the stationary
point. Setting the gradient to zero,
%
\begin{equation}
  -2\,\E[\bm{\theta} \given \mathbf{y}]
  + 2\,\hat{\bm{\theta}} = \mathbf{0}
  \qquad\Longrightarrow\qquad
  \hat{\bm{\theta}}^{*}(\mathbf{y})
  = \E[\bm{\theta} \given \mathbf{y}].
\end{equation}
\end{proof}

\begin{proposition}[Quantile loss]
\label{prop:quantile}
For a scalar parameter $\theta$, under the $\alpha$-quantile loss
$L_\alpha(\theta, \hat{\theta}) = (\alpha -
\mathbf{1}_{\{\theta < \hat{\theta}\}})(\theta - \hat{\theta})$, the
Bayes estimator is the $\alpha$-quantile of the posterior.
\end{proposition}

\begin{proof}
The $\alpha$-quantile loss equals
$(1-\alpha)(\hat{\theta} - \theta)$ when $\theta < \hat{\theta}$ and
$\alpha(\theta - \hat{\theta})$ when $\theta \ge \hat{\theta}$, so the
posterior expected loss splits at $\theta = \hat{\theta}$ as
%
\begin{equation}
  R(\hat{\theta})
  = (1-\alpha)\!\int_{-\infty}^{\hat{\theta}}
    (\hat{\theta} - \theta)\,p(\theta \given \mathbf{y}) \dd\theta
  + \alpha\!\int_{\hat{\theta}}^{\infty}
    (\theta - \hat{\theta})\,p(\theta \given \mathbf{y}) \dd\theta.
\end{equation}
%
Here $\hat{\theta}$ appears both in the limits of integration and in
the integrands, so each term is differentiated by the Leibniz
integral rule. In the first integral the boundary term carries the
integrand evaluated at $\theta = \hat{\theta}$, namely
$(\hat{\theta} - \hat{\theta}) = 0$, so only the integrand derivative
survives:
%
\begin{align}
  \frac{\dd}{\dd\hat{\theta}}
  \int_{-\infty}^{\hat{\theta}}
  (\hat{\theta} - \theta)\,p(\theta \given \mathbf{y}) \dd\theta
  &= \underbrace{(\hat{\theta} - \hat{\theta})\,
     p(\hat{\theta} \given \mathbf{y})}_{=\;0}
   + \int_{-\infty}^{\hat{\theta}}
     \frac{\partial}{\partial\hat{\theta}}
     \big[(\hat{\theta} - \theta)\,
     p(\theta \given \mathbf{y})\big] \dd\theta
   \notag\\[2pt]
  &= \int_{-\infty}^{\hat{\theta}}
     p(\theta \given \mathbf{y}) \dd\theta \\[2pt]
   &= F_{\theta \given \mathbf{y}}(\hat{\theta}),
\end{align}
%
where $F_{\theta \given \mathbf{y}}$ is the posterior cumulative
distribution function. 
In the second integral the boundary term vanishes for the same
reason, and differentiating the lower limit contributes a minus sign:
%
\begin{align}
  \frac{\dd}{\dd\hat{\theta}}
  \int_{\hat{\theta}}^{\infty}
  (\theta - \hat{\theta})\,p(\theta \given \mathbf{y}) \dd\theta
  &= \underbrace{-(\hat{\theta} - \hat{\theta})\,
     p(\hat{\theta} \given \mathbf{y})}_{=\;0}
   - \int_{\hat{\theta}}^{\infty}
     p(\theta \given \mathbf{y}) \dd\theta
   \notag\\[2pt]
  &= -\big(1 - F_{\theta \given \mathbf{y}}(\hat{\theta})\big).
\end{align}
%
Combining the two and simplifying,
%
\begin{align}
  R'(\hat{\theta})
  &= (1-\alpha)\,F_{\theta \given \mathbf{y}}(\hat{\theta})
   - \alpha\big(1 - F_{\theta \given \mathbf{y}}(\hat{\theta})\big)
   \notag\\[2pt]
  &= F_{\theta \given \mathbf{y}}(\hat{\theta}) - \alpha.
\end{align}
%
The posterior expected loss $R$ is convex in $\hat{\theta}$, so its
minimiser is its stationary point. Setting $R'(\hat{\theta}) = 0$
gives
%
\begin{equation}
  F_{\theta \given \mathbf{y}}(\hat{\theta}^{*}) = \alpha,
\end{equation}
%
which is precisely the definition of the $\alpha$-quantile of the
posterior.
\end{proof}

The case $\alpha = \tfrac{1}{2}$ recovers the posterior median under
absolute-error loss; symmetric pairs such as $\alpha \in
\{0.025, 0.975\}$ give the endpoints of a central credible interval.
This thesis works throughout with squared-error loss and hence with
the posterior mean.


% ----------------------------------------------------------------
\section{Neural Bayes estimators}
\label{sec:nbe}

The Bayes estimator~\eqref{eq:pointwise} is a fixed function of the data: it returns the optimal estimate for every possible dataset, not just the one in hand. This is what makes it the right object to approximate: a single approximation, reused across datasets, which amortises the cost of fitting over every later inference, where MCMC must instead be re-run from scratch each time (Section~\ref{sec:mcmc-cost}). Approximating it requires first defining a \emph{parametric family} of estimators in place of the intractable search over all functions, then fitting that family to a Monte Carlo estimate of an otherwise-intractable Bayes risk, and finally using a neural network as a flexible surrogate for this family. The three subsections below develop these in turn.

\paragraph{A parametric family of estimators.}
Definition~\ref{def:bayes-estimator} poses the Bayes estimator as the
solution of a \emph{variational} problem: the minimisation
\eqref{eq:bayes-estimator} ranges over all functions
$\hat{\bm{\theta}} : \mathcal{Y} \to \Theta$, an
infinite-dimensional search that is no more tractable than the
posterior it encodes. We make it tractable by restricting the search
to a family of functions
$\{\hat{\bm{\theta}}_{\bm{\gamma}} : \bm{\gamma} \in \Gamma\}$ indexed
by a finite-dimensional parameter $\bm{\gamma}$. The Bayes risk is then
a function of $\bm{\gamma}$ alone,
%
\begin{equation}
  r(\bm{\gamma})
  = \E_{\bm{\theta}, \mathbf{y}}
    \big[L(\bm{\theta},
    \hat{\bm{\theta}}_{\bm{\gamma}}(\mathbf{y}))\big],
  \label{eq:population-loss}
\end{equation}
%
and the variational problem of Definition~\ref{def:bayes-estimator}
becomes an ordinary parametric optimisation,
%
\begin{equation}
  \bm{\gamma}^{*}
  = \arg\min_{\bm{\gamma} \in \Gamma}\; r(\bm{\gamma}),
  \label{eq:gamma-star}
\end{equation}
%
over the finite-dimensional space $\Gamma$. Restricting to a finite-dimensional search means that the family may not contain the Bayes estimator exactly; the best it can do
is the closest member, and the approximation is only as good as the
family allows.

\paragraph{A Monte Carlo objective.}
The objective~\eqref{eq:population-loss} is still an expectation over
parameter--data pairs, taken with respect to the joint distribution
$p(\bm{\theta})\,p(\mathbf{y} \given \bm{\theta})$, and is no more
available in closed form than the posterior. However, it can be
approximated. The joint distribution is something we can sample
from, even when the likelihood cannot be evaluated: draw a parameter
from the prior, then simulate data from the model given that
parameter. Drawing $K$ such pairs gives the Monte Carlo approximation:
%
\begin{equation}
  r(\bm{\gamma})
  \approx \frac{1}{K}\sum_{k=1}^{K}
  L\big(\bm{\theta}^{(k)},
  \hat{\bm{\theta}}_{\bm{\gamma}}(\mathbf{y}^{(k)})\big),
  \label{eq:empirical-bayes-risk}
\end{equation}
%
the \emph{empirical Bayes risk}, which replaces the intractable
expectation~\eqref{eq:population-loss} by an average we can actually
compute. Minimising~\eqref{eq:empirical-bayes-risk} over $\bm{\gamma}$
is a tractable surrogate for the parametric
problem~\eqref{eq:gamma-star}, and hence for the original problem. The construction is \emph{likelihood-free}: it asks only that
we can simulate from $p(\mathbf{y} \given \bm{\theta})$, never that we
evaluate the likelihood analytically.

\paragraph{A neural network surrogate.}
For the surrogate to be
worthwhile the family $\{\hat{\bm{\theta}}_{\bm{\gamma}}\}$ must be flexible enough to contain a close
approximation to the Bayes estimator, which may be a highly nonlinear
function of the data. Neural networks are the natural choice: by their
universal function-approximation property they can approximate the
estimator map to arbitrary accuracy given sufficient capacity, while
remaining a finite-dimensional family that
\eqref{eq:empirical-bayes-risk} can be optimised over by gradient
methods. Taking $\hat{\bm{\theta}}_{\bm{\gamma}}$ to be a neural
network with weights $\bm{\gamma}$, trained by minimising the empirical
Bayes risk, yields a neural Bayes
estimator~\citep{sainsburydale2024nbe}.

The loss function chosen in~\eqref{eq:empirical-bayes-risk} determines
which posterior summary the trained network approximates, exactly as in
Section~\ref{sec:loss-summaries}. In summary, the procedure followed to obtain the trained network and apply it to an unseen dataset is as follows:

\begin{enumerate}
  \item \textbf{Simulate.} Draw parameters
  $\bm{\theta}^{(k)} \sim p(\bm{\theta})$ from the prior, then
  simulate $n$ data samples
  $\mathbf{y}^{(k)} \sim p(\mathbf{y} \given \bm{\theta}^{(k)})$,
  with $\mathbf{y}^{(k)} \in \reals^{n}$, for $k = 1, \dots, K$.
  \item \textbf{Train.} Choose the weights $\hat{\bm{\gamma}}$ to
  minimise the empirical Bayes risk~\eqref{eq:empirical-bayes-risk}.
  \item \textbf{Apply.} For a new dataset $\mathbf{y}$, report
  $\hat{\bm{\theta}}_{\hat{\bm{\gamma}}}(\mathbf{y})$.
\end{enumerate}

\subsection{Training by stochastic gradient descent}
\label{sec:training-as-mc}

Training minimises the empirical Bayes
risk~\eqref{eq:empirical-bayes-risk} by \emph{stochastic gradient
descent}. Gradient descent is an iterative method: starting from some
initial weights, it repeatedly adjusts $\bm{\gamma}$ by a small step in
the direction of steepest descent (the negative gradient
$-\nabla_{\bm{\gamma}} r(\bm{\gamma})$) so that the risk decreases
from one step to the next, and continues until the risk levels off.
Computing the exact gradient would mean summing over all $K$ simulated
pairs at every step, which is expensive. The stochastic variant
avoids this: at each step it draws a small random \emph{minibatch} of
$B$ pairs and uses the gradient of the loss on that minibatch in place
of the full gradient. Written out, a minibatch
$\{(\bm{\theta}^{(i)}, \mathbf{y}^{(i)})\}_{i=1}^{B}$ of pairs simulated
from the model gives the average of the per-pair gradients,
%
\begin{equation}
  \hat{g}_{B}(\bm{\gamma})
  = \frac{1}{B}\sum_{i=1}^{B}
    \nabla_{\bm{\gamma}}\,
    L\big(\bm{\theta}^{(i)},
    \hat{\bm{\theta}}_{\bm{\gamma}}(\mathbf{y}^{(i)})\big).
  \label{eq:minibatch-gradient}
\end{equation}
%
Each pair is drawn from the joint distribution
$p(\bm{\theta})\,p(\mathbf{y} \given \bm{\theta})$, so interchanging
differentiation and expectation (assuming usual regularity conditions) shows this average to be unbiased for
the gradient of the population risk~\eqref{eq:population-loss},
%
\begin{equation}
  \E\big[\hat{g}_{B}(\bm{\gamma})\big]
  = \nabla_{\bm{\gamma}}\,
    \E_{\bm{\theta}, \mathbf{y}}
    \big[L(\bm{\theta}, \hat{\bm{\theta}}_{\bm{\gamma}}(\mathbf{y}))\big]
  = \nabla_{\bm{\gamma}}\, r(\bm{\gamma}).
  \label{eq:minibatch-unbiased}
\end{equation}
%
The minibatch gradient is therefore an unbiased Monte Carlo estimate of
the population gradient $\nabla_{\bm{\gamma}} r(\bm{\gamma})$, the same
quantity full-batch descent approximates: each step follows the true
gradient on average, cheaply, but with sampling noise.

Like any Monte Carlo estimator, this training gradient has a variance.
The variance is not a property of the inference problem but of how the
loss is sampled, and it can be reduced by the standard Monte Carlo
variance-reduction techniques, applied here to the training, as
distinct from their application to the inference itself, which we
investigate in Chapter~\ref{ch:rao-blackwell} and
Chapter~\ref{ch:importance-sampling}.

% ----------------------------------------------------------------
\section{Neural architectures}
\label{sec:architectures}

The construction of Section~\ref{sec:nbe} leaves the family
$\{\hat{\bm{\theta}}_{\bm{\gamma}}\}$ unspecified beyond the
requirement that it be a neural network. In practice there are actually many different neural network architectures, with different desirable properties useful for solving varying statistical problems. A Bayes estimator inherits the symmetries of the model that
generates the data: if the data are exchangeable the estimator is
invariant to their ordering, if they are a time series the estimator
depends on their order, and so on. Building these properties into the
architecture, rather than hoping the network learns them, reduces the
function class to one that already contains estimators of the right
form; this makes training more efficient and stable. This section
describes three architectures, each matched to a different data
structure: DeepSets for independent and identically distributed (i.i.d.)
observations, recurrent networks for sequential data, and feed-forward networks for data of fixed
dimension with no such structure to exploit.

\subsection{The DeepSets architecture}
\label{sec:deepsets}

Many inference problems involve \emph{replicated} data: a sample
$\mathbf{y} = (y_1, \dots, y_n)$ of independent observations from a
common model. A Bayes estimator for such data is invariant to the
ordering of the observations, and the network should be too. The
DeepSets architecture~\citep{zaheer2017deepsets} enforces this
permutation invariance by construction. Each observation is mapped
through a shared network $\psi$ to a high-dimensional embedding; the embeddings are
combined by a permutation-invariant operation, typically summation;
and a second network $\phi$ maps the aggregate to the estimate,
%
\begin{equation}
  \hat{\bm{\theta}}(\mathbf{y})
  = \phi\!\left(\sum_{i=1}^{n} \psi(y_i)\right).
  \label{eq:deepsets}
\end{equation}
%
The sum $\sum_i \psi(y_i)$ plays the role of a learned summary
statistic. For exponential-family models the sufficient statistic has
exactly this additive form, and the architecture can be read as
learning a sufficient statistic directly from simulated data, without
the likelihood being known. Figure~\ref{fig:deepsets} summarises the
data flow in this architecture.

\begin{figure}[H]
  \centering
  \begin{tikzpicture}[
    >=Stealth,
    obs/.style   = {circle, draw=blue!60, fill=blue!12,
                     minimum size=8mm, font=\small},
    embed/.style = {circle, draw=green!55!black, fill=green!14,
                     minimum size=8mm, font=\small},
    net/.style   = {rectangle, draw=black!70, fill=white,
                     minimum size=7mm, font=\small},
    agg/.style   = {circle, draw=orange!85!black, fill=orange!22,
                     minimum size=9mm, font=\small},
    estim/.style = {circle, draw=red!55, fill=red!12,
                     minimum size=8mm, font=\small},
    lbl/.style   = {font=\footnotesize, text=black!65, align=center},
    every node/.style = {inner sep=2pt},
  ]
  % input observations
  \node[obs] (y1) at (0,3)    {$y_1$};
  \node[obs] (y2) at (0,2)    {$y_2$};
  \node[obs] (y3) at (0,1)    {$y_3$};
  \node      (yd) at (0,0.25) {$\vdots$};
  \node[obs] (yn) at (0,-0.6) {$y_n$};
  % pointwise network psi
  \node[net] (p1) at (2,3)    {$\psi$};
  \node[net] (p2) at (2,2)    {$\psi$};
  \node[net] (p3) at (2,1)    {$\psi$};
  \node[net] (pn) at (2,-0.6) {$\psi$};
  % embeddings
  \node[embed] (s1) at (4,3)    {$s_1$};
  \node[embed] (s2) at (4,2)    {$s_2$};
  \node[embed] (s3) at (4,1)    {$s_3$};
  \node[embed] (sn) at (4,-0.6) {$s_n$};
  % aggregate, summary, final map, estimate
  \node[agg]   (sum) at (6.3,1.2)  {$\textstyle\sum$};
  \node[embed] (S)   at (8.1,1.2)  {$S$};
  \node[net]   (phi) at (9.9,1.2)  {$\phi$};
  \node[estim] (th)  at (11.7,1.2) {$\hat{\bm{\theta}}$};
  % edges
  \foreach \i in {1,2,3,n}{
    \draw[->] (y\i) -- (p\i);
    \draw[->] (p\i) -- (s\i);
    \draw[->] (s\i) -- (sum);
  }
  \draw[->] (sum) -- (S);
  \draw[->] (S)   -- (phi);
  \draw[->] (phi) -- (th);
  % column labels
  \node[lbl] at (0,-1.6)    {Input};
  \node[lbl] at (2,-1.6)    {Pointwise\\network $\psi$};
  \node[lbl] at (4,-1.6)    {Summary\\embeddings};
  \node[lbl] at (6.3,-0.45) {Aggregate};
  \node[lbl] at (9.9,-0.45) {Final map $\phi$};
  \node[lbl] at (11.7,-0.45){Estimate};
  \end{tikzpicture}
  \caption{The DeepSets architecture for permutation-invariant
  estimation. Each observation $y_i$ is mapped by the shared network
  $\psi$ to an embedding $s_i$; the embeddings are summed into the
  permutation-invariant aggregate $S = \sum_i \psi(y_i)$, which the
  second network $\phi$ maps to the estimate
  $\hat{\bm{\theta}}(\mathbf{y})$.}
  \label{fig:deepsets}
\end{figure}








\subsection{Recurrent architectures for sequential data}
\label{sec:recurrent-arch}

DeepSets is appropriate only when the observations are exchangeable.
For \emph{sequential} data such as a time series $\mathbf{y} = (y_1, \dots,
y_n)$ in which the ordering carries information, permutation
invariance would discard precisely the dependence the parameters
govern. The autoregressive model of
Chapter~\ref{ch:importance-sampling} is the case in point: the AR(1)
coefficient is identified by the correlation between successive
observations, which a permutation-invariant summary cannot see.

The natural alternative is a \emph{recurrent} network, which processes
the observations in order while maintaining a hidden state
$\mathbf{h}_t$ that summarises the sequence up to step $t$,
%
\begin{equation}
  \mathbf{h}_t = \psi(\mathbf{h}_{t-1}, y_t),
  \qquad
  \hat{\bm{\theta}}(\mathbf{y}) = \phi(\mathbf{h}_n).
\end{equation}
%
Here $\psi$ is a recurrent cell shared across steps and $\phi$ maps
the terminal hidden state to the estimate; Figure~\ref{fig:recurrent}
shows the network unrolled over the sequence. The hidden state plays
the role that the sum $\sum_i \psi(y_i)$ plays in DeepSets, a learned
summary of the data, but one accumulated in sequence and so sensitive
to temporal order. The experiments of
Chapter~\ref{ch:importance-sampling} take $\psi$ to be a gated
recurrent unit~\citep{cho2014gru}, a widely adopted recurrent block whose gating mitigates the
vanishing gradients over long sequences.


\begin{figure}[H]
  \centering
  \begin{tikzpicture}[
    >=Stealth,
    obs/.style   = {circle, draw=blue!60, fill=blue!12,
                     minimum size=8mm, font=\small},
    net/.style   = {rectangle, draw=black!70, fill=white,
                     minimum size=9mm, font=\small},
    estim/.style = {circle, draw=red!55, fill=red!12,
                     minimum size=8mm, font=\small},
    hid/.style   = {font=\footnotesize, text=green!45!black},
    lbl/.style   = {font=\footnotesize, text=black!65, align=center},
    every node/.style = {inner sep=2pt},
  ]
  % recurrent cells, unrolled over the sequence
  \node[net] (c1) at (0,1)    {$\psi$};
  \node[net] (c2) at (2.4,1)  {$\psi$};
  \node[net] (c3) at (4.8,1)  {$\psi$};
  \node      (cd) at (6.6,1)  {$\cdots$};
  \node[net] (cn) at (8.4,1)  {$\psi$};
  % inputs
  \node[obs] (y1) at (0,-0.7)   {$y_1$};
  \node[obs] (y2) at (2.4,-0.7) {$y_2$};
  \node[obs] (y3) at (4.8,-0.7) {$y_3$};
  \node[obs] (yn) at (8.4,-0.7) {$y_n$};
  % final map and estimate
  \node[net]   (phi) at (10.6,1) {$\phi$};
  \node[estim] (th)  at (12.4,1) {$\hat{\bm{\theta}}$};
  \coordinate (h0) at (-1.7,1);
  % input edges
  \foreach \i in {1,2,3,n}{\draw[->] (y\i) -- (c\i);}
  % recurrent (hidden-state) edges
  \draw[->] (h0) -- node[above,hid]{$\mathbf{h}_0$} (c1);
  \draw[->] (c1) -- node[above,hid]{$\mathbf{h}_1$} (c2);
  \draw[->] (c2) -- node[above,hid]{$\mathbf{h}_2$} (c3);
  \draw[->] (c3) -- node[above,hid]{$\mathbf{h}_3$} (cd);
  \draw[->] (cd) -- node[above,hid]{$\mathbf{h}_{n-1}$} (cn);
  \draw[->] (cn) -- node[above,hid]{$\mathbf{h}_n$} (phi);
  \draw[->] (phi) -- (th);
  % labels
  \node[lbl] at (0,-1.75)   {Input};
  \node[lbl] at (4.8,2.05)  {Recurrent cell $\psi$};
  \node[lbl] at (10.6,0.1)  {Final map $\phi$};
  \node[lbl] at (12.4,0.1)  {Estimate};
  \end{tikzpicture}
  \caption{A recurrent architecture for sequential data, unrolled over
  the sequence. At step $t$ the shared cell $\psi$ combines the
  observation $y_t$ with the previous hidden state $\mathbf{h}_{t-1}$
  to produce $\mathbf{h}_t$; the final map $\phi$ turns the terminal
  hidden state $\mathbf{h}_n$ into the estimate
  $\hat{\bm{\theta}}(\mathbf{y})$.}
  \label{fig:recurrent}
\end{figure}




\subsection{Feed-forward architectures for fixed-dimension data}
\label{sec:mlp-arch}

When the data carry neither exchangeable nor sequential structure to
exploit (as in linear regression with a fixed design matrix
$X$, where each observation is tied to its own row of
covariates and the observations are neither interchangeable nor
ordered in time), there is no symmetry for the architecture to encode.
The natural choice is then a \emph{feed-forward} network, which takes
the data as a fixed-length vector and maps it to the estimate through
a sequence of affine transformations and nonlinearities,
%
\begin{equation}
  \hat{\bm{\theta}}(\mathbf{y}) = \phi(\mathbf{y}).
\end{equation}
%
Unlike DeepSets and the recurrent network, the feed-forward
architecture builds in no invariance and accommodates no change in the
size of the data: the input dimension is fixed by the network, as
suits a problem in which the number of observations and their meaning
are fixed in advance. What it gives up in structure it makes up in
generality, serving as the default when no exploitable structure is
present.






% ----------------------------------------------------------------
\section{Example 1: estimating normal parameters}
\label{sec:normal-example}

As a concrete demonstration, we return to Example~1, the normal model:
$n$ independent observations
$y_i \sim \Normal(\mu, \sigma^2)$ with priors $\mu \sim \Normal(0, 1)$
and $\sigma \sim \Uniform(0, 5)$ and estimand
$\bm{\theta} = (\mu, \sigma)$. The model is a controlled test case
rather than a genuine application: its posterior is available in closed
form up to a normalising constant
(the log-posterior~\eqref{eq:normal-logpost}), so an accurate reference
can be computed for every test dataset and the amortised estimator
checked against it.

That reference is the Metropolis--Hastings sampler of
Section~\ref{sec:cs-normal}. Run separately on each of $256$ test
datasets, its post-burn-in average estimates the posterior mean
$\E[\bm{\theta} \given \mathbf{y}]$ --- the same summary the
squared-error-trained network targets
(Proposition~\ref{prop:squared-error}), so the two methods are compared
on equal terms.

A neural Bayes estimator with the DeepSets
architecture~\eqref{eq:deepsets} was trained under squared-error loss
on simulated data and compared against that reference.
Figure~\ref{fig:normal-model} compares the two, and
Table~\ref{tab:normal-model} summarises their accuracy and cost. They
agree closely in accuracy: the mean squared error of the neural
estimator is within a few percent of that of Metropolis--Hastings for
$\mu$ and somewhat larger for $\sigma$. They differ sharply in cost.
Metropolis--Hastings must be run afresh for each dataset, whereas the
trained network returns an estimate in a single forward pass --- here
roughly $2{,}800$ times faster per dataset.

\begin{figure}[t]
  \centering
  \includegraphics[width=\textwidth]{normal_model_estimates}
  \caption{Estimated against true parameter values for the normal
  model, across $256$ test datasets, for Metropolis--Hastings and
  the neural Bayes estimator. Both methods track the true parameters
  closely; the dashed line marks the hypothetical perfect estimation.}
  \label{fig:normal-model}
\end{figure}

\begin{table}[t]
  \centering
  \caption{Normal model: accuracy (mean squared error over the $256$
  test datasets) and per-dataset cost of the Metropolis--Hastings
  reference and the neural Bayes estimator. The two agree closely in
  accuracy, but the amortised network is faster per dataset by some
  three orders of magnitude (single-threaded CPU timing; $5{,}000$-step
  chains for Metropolis--Hastings).}
  \label{tab:normal-model}
  \begin{tabular}{l rr l}
    \toprule
    & $\mu$ (MSE) & $\sigma$ (MSE) & Cost per dataset \\
    \midrule
    Metropolis--Hastings   & $0.095$ & $0.053$ & $\approx 26$\,ms (re-run per dataset) \\
    Neural Bayes estimator & $0.097$ & $0.075$ & $\approx 9\,\mu$s (one forward pass) \\
    \bottomrule
  \end{tabular}
\end{table}
